\subsection{Real-world validation study}
\label{sec:realworld}

The curated corpus (Section~\ref{sec:evaluation}) was designed to contain
representative defects, so its defect rates measure detection capability rather
than prevalence. To estimate \emph{base rates} on independently authored
workflows, and to measure extractor fidelity on code the authors did not write,
we mined public GitHub repositories.

\paragraph{Mining pipeline.}
Using GitHub code search over the four frameworks, we identified 318 candidate
repositories. We shallow-cloned a 180-repository sample, scanned 383 candidate
files, and extracted 277 workflow graphs from 149 distinct repositories
(LangGraph~206, CrewAI~76).\footnote{AutoGen and ADK repositories ranked later in
the candidate list and fell past the sample cap; balancing all four frameworks is
in progress.} Extraction used the static AST fallback
(\texttt{scripts/ast\_extractor.py}); the runtime extractors of
Section~\ref{sec:graph_model} read the compiled graph object and are expected to
be more faithful, but require per-repository dependency installation that is
infeasible at this scale. Mined workflows are small: a median of 2~non-sentinel
nodes (mean~3.1, max~35), and are predominantly static linear pipelines.

\paragraph{Validation protocol.}
On a stratified 60-workflow sample (40~LangGraph, 20~CrewAI) we ran two
independent LLM-agent passes. First, an agent reconstructed a \emph{ground-truth}
graph from source alone (blind to the extractor output), following the
extractor's node-identifier conventions so the two align; we then computed
node/edge precision--recall and node-kind accuracy against the AST extraction.
Second, for each of the 131 flagged defects, one agent labelled it
\textsc{real}, \textsc{extraction-artifact}, or \textsc{intentional} from source,
and an adversarial verifier independently attempted to overturn each label;
disagreements were recorded as \textsc{arguable}. The verifier overturned only
4.6\% of labels, always toward \textsc{arguable}.

\begin{table}[t]
\centering
\small
\caption{AST-extractor fidelity on real GitHub workflows ($n=60$ ground-truth
graphs). Stub-based validation (Section~\ref{sec:accuracy}) reported perfect
precision/recall; real code is substantially harder, chiefly in edge recovery
and node-kind classification.}
\label{tab:realworld_fidelity}
\begin{tabular}{@{}lrrr@{}}
\toprule
\textbf{Metric} & \textbf{Overall} & \textbf{LangGraph} & \textbf{CrewAI} \\
\midrule
Node precision      & 0.934 & --    & --    \\
Node recall         & 0.896 & --    & --    \\
Edge precision      & 0.871 & 0.956 & 0.700 \\
Edge recall         & 0.685 & 0.682 & 0.692 \\
Node-kind accuracy  & 0.741 & 0.647 & 0.929 \\
\bottomrule
\end{tabular}
\end{table}

\paragraph{Extractor fidelity.}
Table~\ref{tab:realworld_fidelity} reports the result. Node detection remains
strong (P~$=0.93$, R~$=0.90$), but edge recall falls to 0.685 and node-kind
accuracy to 0.741---far below the perfect scores on programmatic stubs. The
dominant kind error is \texttt{tool}$\rightarrow$\texttt{llm} (35 nodes that
invoke tools inside their body without declared bindings), followed by
\texttt{passthrough}$\rightarrow$\texttt{llm} (21) and misclassification of
conditional-edge nodes as routers. Human nodes were misclassified 11~times,
quantifying on real code the naming-heuristic brittleness noted in
Section~\ref{sec:accuracy}.

\begin{table}[t]
\centering
\small
\caption{Adversarial triage of 131 defects flagged on real workflows. Every
structural flag is an extraction artifact; the human-gate policy flag is
predominantly an intentional design choice.}
\label{tab:realworld_triage}
\begin{tabular}{@{}lrrrr@{}}
\toprule
\textbf{Check} & \textbf{Real} & \textbf{Artifact} & \textbf{Intentional} & \textbf{Arguable} \\
\midrule
exit\_reachability     & 0 & 14 & 0  & 0 \\
reverse\_reachability  & 0 & 25 & 0  & 0 \\
dead\_ends             & 0 & 24 & 0  & 1 \\
router\_shape          & 0 & 2  & 0  & 0 \\
tool\_declarations     & 0 & 5  & 0  & 0 \\
human\_presence        & 1 & 3  & 50 & 6 \\
\midrule
\textbf{Total}         & \textbf{1} & \textbf{73} & \textbf{50} & \textbf{7} \\
\bottomrule
\end{tabular}
\end{table}

\paragraph{Defect prevalence.}
Table~\ref{tab:realworld_triage} gives the sobering result: of 131 flagged
defects, only~1 (0.8\%) is a genuine workflow bug. All 65~structural flags are
extraction artifacts---the static checks are only as sound as the recovered
topology, and where the extractor misses edges it reports spurious dead-ends,
unreachable exits, and livelocks. A direct cross-check supports this: mean AST
edge recall was 0.51 for workflows carrying an artifact-labelled defect versus
1.0 for the single real defect. The \texttt{human\_presence} policy check fires
on 98\% of workflows, but 91\% of those firings are intentional (genuinely
low-risk tasks), validating the context-dependence argued in
Section~\ref{sec:evaluation} and motivating a risk-aware variant.

\paragraph{Takeaway.}
Topology-level defects are \emph{rare} in publicly available agent workflows,
which are mostly small static pipelines, and the curated-benchmark rates do not
transfer to production prevalence. The practical determinants of static-analysis
utility are therefore extraction fidelity and policy applicability rather than
the check algorithms themselves. We read this as redirecting the research agenda:
toward higher-fidelity extraction (sandboxed runtime extraction; AST/runtime
cross-validation) and risk-aware policy specification, and away from cataloguing
topology checks whose real-world firing is dominated by extraction noise.
